\section{附录：Lean 与 AI for Math 学习路线}

\noindent\begin{longtable}{p{0.20\textwidth}p{0.35\textwidth}p{0.36\textwidth}}
\toprule
阶段 & 学什么 & 为什么学 \\
\midrule
数学语言 & 命题、定义、lemma、proof、反例 & 理解自然语言证明如何组织。 \\
Lean 基础 & theorem、example、rw、simp、apply、exact、induction & 理解 proof assistant 的交互方式。 \\
Mathlib & 常用定理库、命名、搜索、依赖关系 & AI 证明离不开已有数学库。 \\
Autoformalization & 自然语言到 Lean statement & 连接论文/题目和形式系统。 \\
Proof Search & tactic 组合、goal state、错误修复 & 让模型能在验证器反馈下迭代。 \\
Human Collaboration & 数学家审美、问题选择、解释力 & 判断证明是否有意义，而不只是通过。 \\
\bottomrule
\end{longtable}

\begin{knowledgebox}{学习路线的意义}
这张表把 AI for Math 的学习顺序从“先学模型”改成“先理解数学和形式系统”。对工程读者来说，Lean 基础和 proof search 是进入这一领域的最低门槛；对数学读者来说，autoformalization 和模型反馈是理解 AI 能力边界的入口。
\end{knowledgebox}

\subsection{本章小结}

AI for Math 是跨学科路线：数学语言、形式系统、模型训练和人机协作缺一不可。学习路径越清楚，越不容易把它误解成“让模型多刷题”。

